% ===== 原文 pp.125--129 =====
% [原文 p.125]
% [跨页方框续：橙色「例 5.3：非阿贝尔规范理论」框起于 p.124 页末（由 tr07c 打开该 examplebox，
%  标题「例 5.3：非阿贝尔规范理论」，且不闭合），本文件续写 p.125 的框内正文，
%  并在 (5.69) 之后闭合该 examplebox；按 STYLE_GUIDE §7 跨页方框合并为一个。]
% [Tutor Note: 拼接顺序 tr07c → tr08a 时，这个 examplebox 的开闭标记分处两个文件，
%  这是 p.124/125 交界处唯一的跨文件环境；除此以外本块所有方框都完整落在本块内。]
它有
\begin{equation}
  \widehat{Z}(\mathcal{G}) = \Z_4\,. \tag{5.60}
\end{equation}
此外还有 $2M$ 个复标量，全部按规范群的矢量表示变换，对此有
\begin{equation}
  q^g = 2 \in \Z_4\,, \tag{5.61}
\end{equation}
对每个 $i$ 皆成立。我们调节该理论，使其具有味对称代数
\begin{equation}
  \mathfrak{f} = \mathfrak{sp}(M) = \mathfrak{usp}(2M)\,, \tag{5.62}
\end{equation}
在这些标量按 $2M$ 维基础表示变换的代数下。相应的单连通群为
\begin{equation}
  F = \Sp(M)\,, \tag{5.63}
\end{equation}
对此我们有
\begin{equation}
  \widehat{Z}(F) = \Z_2\,, \tag{5.64}
\end{equation}
而物质场给出荷
\begin{equation}
  q^f = 1 \in \Z_2\,, \tag{5.65}
\end{equation}
于是总共有
\begin{equation}
  Q = \langle (2,1) \rangle \in \Z_4 \times \Z_2\,. \tag{5.66}
\end{equation}
同样有
\begin{equation}
  Q_{\mathrm{genuine}} = Q \cap \widehat{Z}(F) = 0\,, \tag{5.67}
\end{equation}
这意味着
\begin{equation}
  \mathcal{Z} = Z(F) = \Z_2\,, \tag{5.68}
\end{equation}
而 0-形式味对称群为
\begin{equation}
  G^{(0)} = F/\mathcal{Z} = \Sp(M)/\Z_2 = \mathrm{PSp}(M)\,. \tag{5.69}
\end{equation}
\end{examplebox}

% [Tutor Note: 式 (5.60) 可独立核对：$\mathrm{Spin}(4N+2)$ 的中心是 $\Z_4$，而
%  $\mathrm{Spin}(4N)$ 的中心是 $\Z_2\times\Z_2$；原文取 $4N+2$ 正是为了让
%  $\widehat{Z}(\mathcal{G})=\Z_4$（若换成 $4N$，(5.61) 里的 $q^g=2$ 就不再是 $\Z_4$ 的元素）。]
% [Tutor Note: 式 (5.66) 的 $\langle(2,1)\rangle$ 指由单个元素 $(2,1)$ 生成的循环子群，
%  不是 $\langle 2\rangle\times\langle 1\rangle$：$2\cdot(2,1)=(4,2)=(0,0)$，故 $Q\cong\Z_2$；
%  它与 $\widehat{Z}(F)=0\times\Z_2$ 的交只含 $(0,0)$，这正是 (5.67) 说 $Q_{\mathrm{genuine}}=0$ 的理由。]

\subsubsection{Bockstein 2-群对称性}
下面给出 Bockstein 2-群对称性的定义：
\begin{definitionbox}{定义 5.2：Bockstein 2-群对称性}
一个 Bockstein 2-群的 Postnikov 类形如
\begin{equation}
  [\Theta] = \mathrm{Bock}\left(\left[w_2^f\right]\right)\,, \tag{5.70}
\end{equation}
其中 $\left[w_2^f\right]$ 是把背景 $G^{(0)}$ 丛提升为 $F$ 丛的 $\mathcal{Z}$-值障碍类，而 $\mathrm{Bock}$ 是与短正合列
\begin{equation}
  0 \to G^{(1)} \to \mathcal{E} \to \mathcal{Z} \to 0\,, \tag{5.71}
\end{equation}
相关联的 Bockstein 映射，这里 $\mathcal{E}$ 是 $\mathcal{Z}$ 被 $G^{(1)}$ 的某个扩张。关于 Bockstein 映射的更多细节，见 (4.106) 附近的讨论。
\end{definitionbox}

% [原文 p.126]
\para{屏蔽解释。}
正如 1-形式对称群 $G^{(1)}$ 可以用线算符的屏蔽来解释，扩张群 $\mathcal{E}$ 也有同样的解释。

由 §3.1 的讨论回顾：对 $G^{(1)}$，我们计算线算符等价类构成的群 $\mathsf{D}_1$，其等价关系为
\begin{equation}
  L \sim L'\,, \tag{5.72}
\end{equation}
只要线 $L$ 与 $L'$ 之间存在一个局域算符。于是有
\begin{equation}
  \widehat{G}^{(1)} = \mathsf{D}_1\,. \tag{5.73}
\end{equation}
类似地，对 $\mathcal{E}$ 我们计算线算符等价类构成的群 $\widetilde{\mathsf{D}}_1$，其等价关系为
\begin{equation}
  L \sim L'\,, \tag{5.74}
\end{equation}
仅当 $L$ 与 $L'$ 之间存在荷为
\begin{equation}
  q^{\mathcal{Z}} = 0 \in \widehat{\mathcal{Z}}\,. \tag{5.75}
\end{equation}
的局域算符。局域算符的荷 $q^{\mathcal{Z}}$ 由其中心荷
\begin{equation}
  q^f \in \widehat{Z}(F)\,, \tag{5.76}
\end{equation}
通过自然满射
\begin{equation}
  \widehat{Z}(F) \to \widehat{\mathcal{Z}}\,. \tag{5.77}
\end{equation}
得到。注意：尽管真的局域算符都满足 $q^{\mathcal{Z}} = 0$，非真局域算符却未必如此。于是有
\begin{equation}
  \widehat{\mathcal{E}} = \widetilde{\mathsf{D}}_1\,. \tag{5.78}
\end{equation}
显然后一个等价关系比前一个更细，因此存在满同态
\begin{equation}
  \widetilde{\mathsf{D}}_1 \to \mathsf{D}_1\,, \tag{5.79}
\end{equation}

% [Tutor Note: (5.74) 与 (5.72) 写成同一个记号 $L\sim L'$，区别只在附加条件 (5.75)；
%  (5.79) 的满射性正是「更细的等价关系给出更少的等价类」：$\widetilde{\mathsf{D}}_1$ 的每个类
%  在 $\mathsf{D}_1$ 中恰好落入一个类，且映射是满的。]

% [原文 p.127]
或等价地
\begin{equation}
  \widehat{\mathcal{E}} \to \widehat{G}^{(1)}\,, \tag{5.80}
\end{equation}
它可以补全为短正合列
\begin{equation}
  0 \to \widehat{\mathcal{Z}} \to \widehat{\mathcal{E}} \to \widehat{G}^{(1)} \to 0\,, \tag{5.81}
\end{equation}
它是 (5.71) 的 Pontryagin 对偶。这就用上述屏蔽解释重新得到了 (5.71)。

\para{2-群背景场。}
至此我们已习惯于：形如
\begin{equation}
  \delta B_2 + \mathrm{Bock}\left(w_2^f\right) = 0\,, \tag{5.82}
\end{equation}
的关系意味着应当把 1-形式背景 $B_2$ 与障碍 $w_2^f$ 合并成单个 $\mathcal{E}$-值闭背景场
\begin{equation}
  B_2^w = i(B_2) + \widetilde{w}_2^f\,, \tag{5.83}
\end{equation}
其中 $i$ 是 (5.71) 中的单同态，而 $\widetilde{w}_2^f$ 是利用 (5.71) 中的满同态对 $w_2^f$ 所作的 $\mathcal{E}$-值提升。

\para{规范理论中的 Bockstein 2-群。}
Bockstein 2-群对称性在规范理论中无处不在。我们从 §5.3.1 所考虑的那类规范理论出发（那里已算出 0-形式对称群），并把电 1-形式对称性一并纳入其中。

我们希望计算 $\widetilde{\mathsf{D}}_1$，它容易辨认为
\begin{equation}
  \widehat{\mathcal{E}} = \widetilde{\mathsf{D}}_1 = \frac{\widehat{Z}(\mathcal{G}) \times \widehat{Z}(F)}{Q}\,, \tag{5.84}
\end{equation}
其中 $Q$ 是 (5.47) 附近描述的荷矩阵。于是有
\begin{equation}
  \mathcal{E} = \text{$Z(\mathcal{G}) \times Z(F)$ 中使 $Q$ 不变的元素构成的子群}\,. \tag{5.85}
\end{equation}
几乎按定义即可看出，$G^{(1)}$ 就是 $\mathcal{E}$ 中使 $Q$ 不变、且投影到 $Z_F$ 上为平凡的元素所构成的子群
\begin{equation}
  G^{(1)} = \{(*,0) \in \mathcal{E} \subseteq Z(\mathcal{G}) \times Z(F)\}\,. \tag{5.86}
\end{equation}
这就可以补全为决定 Bockstein 2-群对称性的短正合列 (5.71)。

% [Tutor Note: (5.84) 与 (5.85) 互为 Pontryagin 对偶，可拿例 5.2 验算：
%  $\widehat{Z}(\mathcal{G})=\Z$、$\widehat{Z}(F)=\Z_2$、$Q=\langle(2,1)\rangle$，商
%  分子含无限群 $\Z$，不能按有限群阶除法 $8/2=4$ 处理；改由 (5.88) 直接读出对偶侧的商正是 (5.87) 的 $\Z_4$。
%  注意分子上带帽、分母 $Q$ 不带帽，这是对偶侧的写法，不要写成 $Z(\mathcal{G})\times Z(F)/Q$。]

% [原文 p.128]
\begin{examplebox}{例 5.4：阿贝尔规范理论}
我们回到例 5.2，计算该阿贝尔规范理论中存在的 Bockstein 2-群对称性。我们有
\begin{equation}
  \widehat{\mathcal{E}} = \frac{\Z \times \Z_2}{\langle (2,1)\rangle} = \Z_4\,, \tag{5.87}
\end{equation}
这容易看出，因为
\begin{equation}
  (0,1) \sim (-2,0)\,,\qquad (4,0) \sim (0,0)\,. \tag{5.88}
\end{equation}
短正合列 (5.71) 因此必为
\begin{equation}
  0 \to \Z_2 \to \Z_4 \to \Z_2 \to 0\,. \tag{5.89}
\end{equation}
的形式。这样的短正合列是唯一的，所以我们无需再算别的。2-群关系形如
\begin{equation}
  \delta B_2 = \mathrm{Bock}\left(w_2^f\right) = w_3^f\,, \tag{5.90}
\end{equation}
其中 $w_2^f$ 与 $w_3^f$ 是
\begin{equation}
  G^{(0)} = \SO(3)\,, \tag{5.91}
\end{equation}
丛的第二与第三 Stiefel--Whitney 类；这里用到了这样一个性质：与短正合列 (5.89) 相关联的 Bockstein 映射把第二 Stiefel--Whitney 类送到第三 Stiefel--Whitney 类。
\end{examplebox}

\begin{examplebox}{例 5.5：非阿贝尔规范理论}
我们回到例 5.3，计算该非阿贝尔规范理论中存在的 Bockstein 2-群对称性。我们有
\begin{equation}
  \widehat{\mathcal{E}} = \frac{\Z_4 \times \Z_2}{\langle (2,1)\rangle} = \Z_4\,, \tag{5.92}
\end{equation}
而相关的短正合列同样是 (5.89)，这意味着存在形如
\begin{equation}
  \delta B_2 = \mathrm{Bock}\left(w_2^f\right)\,, \tag{5.93}
\end{equation}
的 2-群对称性，其中 $w_2^f$ 是把 $G^{(0)} = \mathrm{PSp}(M)$ 丛提升为 $F = \Sp(M)$ 丛的障碍。
\end{examplebox}

% [Tutor Note: (5.87)--(5.88) 可用群论核实：由 $(2,1)=0$ 与 $(0,2)=0$（$\Z_2$ 分量）得
%  $(4,0)=2(2,0)=2[(2,1)-(0,1)]=2(0,1)=(0,0)$，故 $(1,0)$ 的阶为 4，商群为 $\Z_4$；
%  (5.88) 的 $(0,1)\sim(-2,0)$ 只是把 $(0,1)$ 用 $(2,1)$ 平移。
%  (5.89) 的唯一性：$\mathrm{Ext}^1(\Z_2,\Z_2)=\Z_2$ 给出两种扩张，分裂的 $\Z_2\times\Z_2$ 与
%  不分裂的 $\Z_4$，原文取的是后者。]

\para{混合规范-味丛。}
回顾：规范理论中电 1-形式对称性的背景场 $B_2$，描述的是把 $\mathcal{G}/G^{(1)}$ 丛提升为 $\mathcal{G}$ 丛的障碍，

% [原文 p.129]
而规范理论要对这些丛求和。它现在不再闭合，意味着该规范理论不再严格地对 $\mathcal{G}/G^{(1)}$ 丛求和。此时应考虑的正确的李群是所谓的\textbf{结构群}（structure group）
\begin{equation}
  \mathcal{S} = \frac{\mathcal{G} \times F}{\mathcal{E}}\,. \tag{5.94}
\end{equation}
一个 0-形式背景固定了 0-形式群 $G^{(0)}$ 的一个丛 $B_1$。规范理论对 $\mathcal{S}$ 的丛求和，这些丛投影到 $B_1$，且提升为 $\mathcal{G} \times F$ 丛的障碍由 2-群背景场 $B_2^w$ 给出。

% [Tutor Note: (5.94) 的结构群 $\mathcal{S}$ 是 0-形式味对称性与电 1-形式对称性的「共同」丛结构：
%  取 $\mathcal{E}$ 为 (5.85) 的子群，则 $\mathcal{S}=(\mathcal{G}\times F)/\mathcal{E}$ 同时映到
%  $\mathcal{G}/G^{(1)}$ 与 $G^{(0)}=F/Z(F)$，这正是混合规范-味丛的完整框架；
%  当 $w_2^f=0$ 时 $\mathcal{E}=G^{(1)}\times 1$，于是 $\mathcal{S}\to\mathcal{G}\times F$ 可裂。]

\section{高阶形式对称性与弦论}
本节简要说明，在量子场论的弦论构造中理解高阶形式对称性所涉及的若干关键想法。我们将考虑两类构造。第一类是把弦论/M-理论紧致化并取解耦动力学引力的极限。第二类则是弦论/M-理论中的全息构造。

\subsection{来自几何工程的高阶形式对称性}
许多有意思的量子场论是强耦合的，没有常规的拉氏量描述。特别是 5d 与 6d 的超共形场论（SCFT），以及 4d 的 Argyres--Douglas 理论。这类理论可以通过把超弦理论或 M-理论紧致化到内禀紧致流形上来构造。这种构造被称为对所研究 QFT 的\textbf{几何工程}（geometric engineering）。这里我们说明 p-形式对称群如何编码在几何工程构造中，主要依据文献 [62--64]。

\para{构造。}
考虑一个 $\mathcal{D}$ 维的弦论/M-理论，其中
\begin{equation}
  \mathcal{D} \in \{10, 11\}\,, \tag{6.1}
\end{equation}
紧致化在形如
\begin{equation}
  M_{\mathcal{D}} = M_d \times M_D\,, \tag{6.2}
\end{equation}
的流形上，其中 $M_d$ 是被几何工程化的量子场论 $\mathcal{T}$ 所处的时空，而 $M_D$ 是内禀紧致化流形。关键是，我们假设内禀流形非紧，即它在无穷远处有边界。几何工程中众所周知，紧致化理论的牛顿常数 $G_N$ 由紧致化流形 $M_D$ 的体积 $\vol(M_D)$ 的倒数设定，而为了构造

% [Tutor Note: (6.1)--(6.2) 中维数用花体 $\mathcal{D}$，而内部流形用普通斜体 $M_D$、
%  $d$ 维时空用 $M_d$：三者不可混同。原文在 §6 后文（如 (6.20)、(6.29)）把维数改回普通
%  斜体 $D$，译文中照原文分别保留 $\mathcal{D}$ 与 $D$ 的写法。]
